% Spanish – Midterm Study Guide (Units 1–3) — EU business casual
% Compile with: pdflatex study-guide.tex

\documentclass[8pt,landscape,a4paper]{article}
\usepackage[utf8]{inputenc}
\usepackage[T1]{fontenc}
\usepackage[left=0.8cm,right=2.6cm,top=0.8cm,bottom=0.8cm,marginparwidth=2cm,marginparsep=0.3em]{geometry}
\usepackage{multicol}
\usepackage{marginnote}
\usepackage{enumitem}
\setlist{nosep,leftmargin=*}
\usepackage{booktabs}
\usepackage{parskip}
\usepackage{xcolor}
\usepackage{fancyhdr}
% EU business casual palette: navy accent, charcoal, soft grey
\definecolor{eublue}{RGB}{0,51,153}
\definecolor{charcoal}{RGB}{51,51,51}
\definecolor{softgrey}{RGB}{240,240,240}
\definecolor{notegrey}{RGB}{90,90,90}
\newcommand{\handnote}[1]{\marginnote{\raggedright\footnotesize\color{notegrey}\textit{#1}\vspace{0.4em}}}
\newcommand{\hlpink}[1]{\textcolor{charcoal}{\textbf{#1}}}
\newcommand{\hlpeach}[1]{\textcolor{eublue}{\textbf{#1}}}
\setlength{\parskip}{0.2em}
\setlength{\columnsep}{0.5em}
\pagestyle{fancy}
\fancyhf{}
\fancyfoot[L]{\raisebox{0.8em}{\small\color{notegrey} Spanish for Beginners}}
\fancyfoot[R]{\raisebox{0.8em}{\small\color{notegrey}\thepage}}
\fancypagestyle{plain}{\fancyhf{}\fancyfoot[L]{\raisebox{0.8em}{\small\color{notegrey} Spanish for Beginners}}\fancyfoot[R]{\raisebox{0.8em}{\small\color{notegrey}\thepage}}}
\raggedright

\begin{document}
\begin{center}
\vspace{0.2em}
{\normalsize\bfseries\color{eublue} Spanish for Beginners -- Midterm Study Guide (Units 1--3)}\\[0.2em]
\footnotesize\textcolor{notegrey}{Use for quick recall; mock exams: datos personales, por/para, nacionalidades, pronunciaci\'on.}\\[0.35em]
\noindent\colorbox{softgrey}{\parbox{0.92\textwidth}{\centering\footnotesize\color{charcoal}
\textbf{Key:} \textbf{Definitions} \quad \textcolor{eublue}{\textbf{Frameworks}} \quad \textit{Margin notes}
\vspace{0.25em}}}
\end{center}
\vspace{-0.4em}

\begin{multicols}{2}

\section*{\color{eublue}Unit 1 -- Preguntas de identificaci\'on}\handnote{Match question to answer.}
\begin{itemize}[leftmargin=*,nosep]
\item \textbf{¿C\'omo te llamas?} $\to$ Francisco. / Me llamo...
\item \textbf{¿Cu\'antos a\~nos tienes?} $\to$ Treinta y cinco.
\item \textbf{¿Cu\'al es tu apellido?} $\to$ Mart\'inez.
\item \textbf{¿De d\'onde eres?} $\to$ De Valencia.
\item \textbf{¿En qu\'e trabajas?} $\to$ Soy camarero.
\end{itemize}

\section*{\color{eublue}Ser, tener, llamarse}\handnote{Conjugations for descriptions.}
\begin{itemize}[leftmargin=*,nosep]
\item \textbf{Ser}: soy, eres, es, somos, sois, son. (identity, origin, profession)
\item \textbf{Tener}: tengo, tienes, tiene, tenemos, ten\'eis, tienen. (age: tengo 20 a\~nos)
\item \textbf{Llamarse}: me llamo, te llamas, se llama... (name)
\item \textit{Mi amiga se llama Sandra y es de Madrid. Las dos somos estudiantes. Sandra tiene veinte a\~nos.}
\end{itemize}

\subsection*{\hlpeach{Ser \& Tener conjugations}}
\begin{center}
\footnotesize
\begin{tabular}{@{}lll@{}}
\toprule
\textbf{Persona} & \textbf{Ser} & \textbf{Tener} \\
\midrule
yo & soy & tengo \\
t\'u & eres & tienes \\
\'el/ella/usted & es & tiene \\
nosotros & somos & tenemos \\
vosotros & sois & ten\'eis \\
ellos/ustedes & son & tienen \\
\bottomrule
\end{tabular}
\end{center}

\section*{\color{eublue}N\'umeros (0--50)}\handnote{Escribe con letras.}
\begin{itemize}[leftmargin=*,nosep]
\item 0--15: cero, uno, dos, tres, cuatro, cinco, \textbf{seis}, siete, ocho, nueve, diez, once, doce, trece, catorce, quince.
\item 16--19: diecis\'eis, diecisiete, dieciocho, diecinueve.
\item 20, 30, 40, 50: veinte, treinta, cuarenta, cincuenta.
\end{itemize}\handnote{6 = seis (not ``sies'').}

\section*{\color{eublue}Nacionalidades \& profesiones}\handnote{Agreement: -o (m), -a (f).}
\begin{itemize}[leftmargin=*,nosep]
\item \textbf{Terminaciones}: argentino/a, portugu\'es/a, belga, canadiense, costarricense, estadounidense (invariable).
\item \textbf{Profesiones}: periodista, arquitecto/a, estudiante, profesor/a, dise\~nador/a, secretaria, cient\'ifica, camarero/a, cocinero/a.
\item \textbf{-ista, -ense, -ante}: same m/f (periodista, canadiense, estudiante).
\end{itemize}

\subsection*{M\'as nacionalidades}
\begin{itemize}[leftmargin=*,nosep]
\item Colombia $\to$ colombiano/a. Ecuador $\to$ ecuatoriano/a. Venezuela $\to$ venezolano/a.
\item Honduras $\to$ hondure\~no/a. Italia $\to$ italiano/a. Holanda $\to$ holand\'es/a.
\item Alemania $\to$ alem\'an/alemana. Francia $\to$ franc\'es/francesa.
\end{itemize}

\subsection*{\hlpink{Lenguas oficiales de Espa\~na}}
\begin{itemize}[leftmargin=*,nosep]
\item \textbf{5 lenguas}: espa\~nol, euskera, catal\'an, aran\'es, gallego.
\item Rom\'anicas (del lat\'in): espa\~nol, catal\'an, gallego, aran\'es. Euskera: origen incierto.
\item Euskera: Pa\'is Vasco, Navarra. Catal\'an: Catalu\~na, Valencia (valenciano), Baleares. Aran\'es: Val d'Aran.
\end{itemize}

\section*{\color{eublue}Lenguas por pa\'is}\handnote{Gentilicio = lengua.}
\begin{itemize}[leftmargin=*,nosep]
\item Rusia $\to$ ruso. Italia $\to$ italiano. Rep.\ Checa $\to$ checo. China $\to$ chino. Brasil $\to$ portugu\'es.
\item Francia $\to$ franc\'es. Alemania $\to$ alem\'an. Jap\'on $\to$ japon\'es. Corea $\to$ coreano.
\end{itemize}

\section*{\color{eublue}Actividades para aprender}\handnote{Ver, escuchar, leer, escribir, practicar, hacer.}
\begin{itemize}[leftmargin=*,nosep]
\item \textbf{ver}: pel\'iculas, series, la televisi\'on, v\'ideos.
\item \textbf{escuchar}: m\'usica, la radio, podcasts.
\item \textbf{leer}: libros, novelas, art\'iculos.
\item \textbf{escribir}: un diario, mensajes, emails.
\item \textbf{practicar}: con nativos, la pronunciaci\'on, la gram\'atica.
\item \textbf{hacer}: ejercicios, deberes, un intercambio.
\end{itemize}

\section*{\color{eublue}Por vs.\ Para}\handnote{Por = reason. Para = purpose.}
\begin{itemize}[leftmargin=*,nosep]
\item \textbf{Para} + infinitivo: purpose. \textit{Quiero aprender espa\~nol para hablar con mis amigos.}
\item \textbf{Porque} + clause: reason. \textit{Quiero visitar el Prado porque quiero ver los cuadros.}
\item \textbf{Por} + sustantivo: reason. \textit{Quiero visitar Roma por la comida y los monumentos.}
\item \textbf{Trick}: Para + infinitive. Porque + conjugated verb. Por + noun.
\item \textbf{M\'as ejemplos}: vivir con familia \textbf{para} practicar; ir a Cuba \textbf{para} visitar La Habana; ir a conciertos \textbf{porque} quiero conocer grupos.
\end{itemize}\handnote{Para = in order to. Porque = because. Por = for/because of.}

\section*{\color{eublue}C o qu?}\handnote{Before e, i: use qu.}
\begin{itemize}[leftmargin=*,nosep]
\item \textbf{qu} before e, i: qui\'en, qu\'e, quince, arquitectura, marroqu\'i.
\item \textbf{c} elsewhere: once, camarero, cocinero, cient\'ifica, gracias, comercial, cu\'antos, cuatro.
\item \textbf{Exception}: \textbf{qu\'ien} needs \textbf{qu} (not just q).
\end{itemize}

\subsection*{Palabras C o qu (lista)}
\begin{itemize}[leftmargin=*,nosep]
\item \textbf{qu}: qui\'en, qu\'e, quince, arquitectura, marroqu\'i, \textbf{qu\'ien}.
\item \textbf{c}: once, camarero, cocinero, cient\'ifica, gracias, comercial, cu\'antos, cuatro.
\end{itemize}

\section*{\color{eublue}Pronunciaci\'on -- Carlos, Garc\'ia, Javier, Celia}\handnote{Sound reference words.}
\begin{itemize}[leftmargin=*,nosep]
\item \textbf{como Carlos}: C before a, o, u; Q = [k]. (colombiana, Costa, Quito, ecuatoriano)
\item \textbf{como Garc\'ia}: G before a, o, u = [g]. (Bogot\'a)
\item \textbf{como Javier}: J; G before e, i = [x]. (argentina, San Jos\'e)
\item \textbf{como Celia}: C before e, i; Z = [s] or [$\theta$]. (Rica, Tegucigalpa, Asunci\'on)
\end{itemize}

\section*{\color{eublue}Listening -- Datos personales}\handnote{5 correct: name, origin, nationality, age, profession.}
\begin{itemize}[leftmargin=*,nosep]
\item Mark: nombre, lugar de origen, nacionalidad, edad, profesi\'on.
\item Not: lengua materna, aficiones, lenguas que aprende (unless said).
\end{itemize}

\section*{\color{eublue}Listening -- Una fiesta}\handnote{True/False. Listen carefully.}
\begin{itemize}[leftmargin=*,nosep]
\item Ana: not from Argentina; not cocinera. Elvira: colombiana. Stephanie: not Dutch. Jos\'e Antonio: profesor de espa\~nol.
\end{itemize}

\section*{\color{eublue}Presente -- verbos regulares}\handnote{-ar, -er, -ir.}
\begin{itemize}[leftmargin=*,nosep]
\item \textbf{-ar}: hablo, hablas, habla, hablamos, habl\'ais, hablan.
\item \textbf{-er}: como, comes, come, comemos, com\'eis, comen.
\item \textbf{-ir}: vivo, vives, vive, vivimos, viv\'is, viven.
\item \textbf{Irregulares}: ser, estar, ir, tener, hacer.
\end{itemize}

\section*{\color{eublue}Vocabulario \'util}\handnote{Palabras clave para ejercicios.}
\begin{itemize}[leftmargin=*,nosep]
\item \textbf{Lugares}: Madrid, Valencia, Bogot\'a, Buenos Aires, San Jos\'e, Quito, Tegucigalpa, Asunci\'on, Bruselas, Oporto, Mil\'an.
\item \textbf{Datos}: nombre, apellido, edad, profesi\'on, nacionalidad, lugar de origen.
\item \textbf{Conectores}: y (and), pero (but), porque (because), de (from/of).
\end{itemize}

\section*{\color{eublue}Quick revision}\handnote{Run through before exam.}
\begin{itemize}[leftmargin=*,nosep]
\item Preguntas de identificaci\'on. Ser, tener, llamarse.
\item N\'umeros 0--50. Nacionalidades \& profesiones (terminaciones).
\item 5 lenguas de Espa\~na. Actividades (ver, escuchar, etc.).
\item Por vs.\ para vs.\ porque. C o qu. Pronunciaci\'on (Carlos, Garc\'ia, Javier, Celia).
\end{itemize}\handnote{If you can explain all these, you're ready.}

\section*{\hlpeach{Exam tips}}\handnote{Read \& listen carefully.}
\begin{itemize}[leftmargin=*,nosep]
\item \textbf{Agreement}: Nacionalidad \& profesi\'on match person (m/f).
\item \textbf{6} = seis. \textbf{qui\'en} = qu (not q).
\item \textbf{Listening}: Focus on name, origin, nationality, profession, age.
\item \textbf{Gapfill}: Read full sentence for context. Check gender (mi amiga $\to$ -a).
\item \textbf{Por/para}: Infinitive after? $\to$ para. Noun after? $\to$ por. Clause with verb? $\to$ porque.
\end{itemize}

\end{multicols}

\newpage
\vspace*{0.3em}
\begin{center}
\textbf{\color{eublue}Quick recap -- last page}\\[0.3em]
\small\textcolor{notegrey}{(Section summaries, not margin notes)}
\end{center}
\vspace{0.4em}
\noindent\color{notegrey}
\setlength{\parskip}{0.3em}
\begin{multicols}{2}

\textbf{Preguntas}\\
¿C\'omo te llamas? ¿Cu\'antos a\~nos? ¿De d\'onde? ¿En qu\'e trabajas? ¿Cu\'al es tu apellido?

\textbf{Ser, tener, llamarse}\\
Soy, tengo, me llamo. Agreement.

\textbf{N\'umeros}\\
0--15, 16--19, 20--50. 6 = seis.

\textbf{Nacionalidades \& profesiones}\\
-o/-a. -ista, -ense invariable.

\textbf{5 lenguas Espa\~na}\\
Espa\~nol, euskera, catal\'an, aran\'es, gallego.

\textbf{Por/Para/Porque}\\
Para + inf. Porque + verb. Por + noun.

\textbf{C o qu}\\
qu before e, i. qui\'en = qu.

\textbf{Pronunciaci\'on}\\
Carlos [k], Garc\'ia [g], Javier [x], Celia [s].

\textbf{Listening}\\
Datos: nombre, origen, nacionalidad, edad, profesi\'on.

\textbf{Presente}\\
-ar, -er, -ir. Ser, estar, tener.

\textbf{Vocabulario}\\
Lugares, datos, conectores.

\end{multicols}
\vfill
\normalsize\normalfont\color{black}
\end{document}
